% Generated draft: requires mathematical and source review.
\documentclass[11pt]{article}
\usepackage[margin=25mm]{geometry}
\usepackage{fontspec}
\setmainfont{Latin Modern Roman}
\usepackage{amsmath,amssymb,mathtools,unicode-math}
\setmathfont{Latin Modern Math}
\usepackage{enumitem,needspace,xurl,hyperref}
\setlength{\parindent}{0pt}
\setlength{\parskip}{5pt}
\setlength{\emergencystretch}{3em}
\newcommand{\N}{\mathbb{N}}
\newcommand{\Z}{\mathbb{Z}}
\newcommand{\Q}{\mathbb{Q}}
\newcommand{\R}{\mathbb{R}}
\newcommand{\C}{\mathbb{C}}
\newcommand{\F}{\mathbb{F}}
\newcommand{\A}{\mathbb{A}}
\newcommand{\PP}{\mathbb{P}}
\newcommand{\E}{\mathbb{E}}
\newcommand{\eps}{\varepsilon}
\newcommand{\dd}{\,\mathrm{d}}
\begin{document}
\section*{UnsolvedMath ID: 700299}
\textbf{Source problem code:} catalogue entry
\subsection{Definitions and mathematical statement}
Let \(S\) be a finite set, let \(F(S)\) be the free group on \(S\), and let \(R\subseteq F(S)\) be finite. The presentation \(\langle S\mid R\rangle\) defines the finitely presented group \(G=F(S)/\langle\!\langle R\rangle\!\rangle\), where \(\langle\!\langle R\rangle\!\rangle\) is the normal closure of \(R\). The group \(G\) is residually finite if, for every \(g\in G\setminus\{1_G\}\), there exist a finite group \(H\) and a homomorphism \(\phi:G\to H\) such that \(\phi(g)\ne 1_H\).

Fix a standard effective encoding of finite words over \(S\cup S^{-1}\) by natural numbers, with primitive recursive syntactic operations and validity test. Define the word-problem characteristic function by
\[
\chi_{G,S}(n)=
\begin{cases}
1,&\text{if }n\text{ encodes a word representing }1_G,\\
0,&\text{otherwise, including invalid codes.}
\end{cases}
\]
A recursive solution of the word problem means that \(\chi_{G,S}\) is a total recursive function. A primitive recursive solution means that \(\chi_{G,S}\) is primitive recursive.

Does there exist a finite presentation \(\langle S\mid R\rangle\) whose group \(G\) is residually finite and for which \(\chi_{G,S}\) is recursive but not primitive recursive? Equivalently, the requested decision function must be computable by an algorithm that halts on every input, while
\[
\forall f:\N\to\{0,1\}\;\bigl(f\text{ primitive recursive}\Longrightarrow \exists n\in\N\; f(n)\ne\chi_{G,S}(n)\bigr).
\]
The supplied statement does not specify a word encoding or formally define a primitive recursive solution. The formulation above uses the decision-function interpretation with standard primitive recursive coding; this interpretive qualification requires review against the original problem.
\subsection{Short English statement}
Does some finitely presented residually finite group have a word problem that an algorithm can decide on every input, but that no primitive recursive decision function can decide? The supplied record labels this problem open; that label is not an independently verified assessment.
\subsection{Sources}
\begin{itemize}
\item Local UnsolvedMath registry, permanent ID \(700299\), exact source problem code \texttt{catalogue entry}, title \emph{Kourovka Notebook Problem 5.15}. The supplied record attributes the proposal to F. B. Cannonito in 1976 and labels its status \texttt{open}; the record was updated on 2026-10-16. These are supplied registry assertions, not independent verification.
\item \emph{The Kourovka Notebook (Unsolved Problems in Group Theory)}, original problem 5.15, 5th issue (1976); the supplied background identifies the 21st issue edition (2026), arXiv:1401.0300v48. The supplied collection description names E. I. Khukhro and V. D. Mazurov as editors and describes coverage of Issues 1--21 (1965--2026).
\end{itemize}
\end{document}
